\documentclass[11pt]{article}
\usepackage{amsmath,longtable,geometry,hyperref}
\geometry{margin=1in}
\title{Optimal Aerodynamic Shape Design for Electric Vehicles}
\author{S-wind Research Team}
\date{January 2026}
\begin{document}
\maketitle

\begin{abstract}
Exhaustive search of 2,935,296 configurations for optimal aerodynamic shape design for electric vehicles. Complete 7-step derivation chain.
\end{abstract}

\section{Problem}
\textbf{Objective}: Minimize drag coefficient C\_D subject to volume efficiency > 0.7

\section{Derivation Chain}
\begin{longtable}{|c|p{2cm}|p{2cm}|p{3cm}|p{2cm}|p{3cm}|}
\hline
Step & Desc & In & Out & Rule & Why\\\hline\endhead
1 & Frontal area & $W, H$ & $A = 0.85×W×H$ & Empirical & Vehicle cross-section ~85\% of rectangle\\\hline
2 & Base drag & $$ & $C_D0 = 0.22$ & Initial & Modern EV baseline\\\hline
3 & Windshield effect & $α_w$ & $ΔC_D,w = 0.002|α_w-30|$ & Semi-empirical & Optimal angle 30°\\\hline
4 & Rear effect & $α_r$ & $ΔC_D,r = 0.003|α_r-15|$ & Semi-empirical & Optimal fastback 15°\\\hline
5 & Total drag & $C_D0, ΔC_D$ & $C_D = C_D0 + Σ ΔC_D$ & Superposition & Sum of components\\\hline
6 & Aero power & $C_D, A, v$ & $P = ½ρC_DAv³$ & First principles & Power = Drag × Velocity\\\hline
7 & Range & $E, P, v$ & $R = E×η/P×v$ & Energy balance & Range = Energy/Power×Speed\\\hline
\end{longtable}

\section{Results}
Optimal objective: 8.5680

\end{document}
